H. R. 9510, the Verification Before Generation in Physical AI Oncology Trials Act of
2026, would amend the Federal Food, Drug, and Cosmetic Act to require that an
automated verification, validation, and uncertainty quantification examination clear
robot-patient interaction code before that code is generated or executed in an
oncology clinical investigation. This paper is a passage framework: it organizes the
case for the bill as the eight questions a member of Congress asks before a yes vote, pairs with the corresponding National
Platform, and the public record. The
eight questions are the mandate for a statute, Congress's authority to act, the
safety the bill creates, the fiscal score, the benefit to constituents,
bipartisanship, the supporting coalition, and the passage path itself. The answers
are concrete: a ten-gate safety floor bound to published consensus standards; an
authorization of 58 million dollars over five years with no new mandatory spending
and a verification cost roughly nineteen times lower than conventional review; and a
platform that, in validated simulation, treated 168 patients with 29 robots at 99.7
percent uptime with zero patient harm events. The framework is written for the
staffer, the member, and the committee counsel who must decide whether to mark up,
cosponsor, and vote for the bill in both chambers. It is a new, analogous
companion with one objective: making H. R. 9510 most likely to pass in both the
House and the Senate.
